% ch11_a 章首+§11.1–11.2 (书页 375–385 / p389–p399)
% 第11章 SUPERCONDUCTIVITY
% 页码核对说明：经与扫描页核对，第 11 章实际始于书页 379（p393.png）；
% p389–p392（书页 375–378）为第 10 章 §10.12 的结尾，归第 10 章文件。
% 本文件实际覆盖书页 379–385（p393–p399）：章首 + §11.1 + §11.2。

\chapter{超导电性}

\epigraphbook{既然变化本是自然的法则，\\ 唯独恒常反而显得奇怪。}{EARL OF ROCHESTER}

\section{电子间的吸引}

超导电性（superconductivity）长期以来被看作金属诸性质中最异乎寻常、最神秘的一种；但是 Bardeen、Cooper 与 Schrieffer 的理论——即\emph{BCS 理论}——已经解释了如此之多，以致我们可以说，如今我们对超导态的理解，几乎不亚于对“正常”态的理解。我们并不打算在这里论及这个庞大而进展迅速的领域的全部内容；重点将放在那些原子过程上，许多不寻常的宏观现象正是由它们引发的。

%FIG{199}: {电子附近晶格的极化。}

整个效应源自任意两个能量几乎相同的电子之间的一种微弱吸引力。我们通常认为自由电子通过其库仑相互作用而相互\emph{排斥}；不过，正如 §5.8 所示，在长距离上这个场会因“其他”电子的屏蔽（screening）效应而显著减弱。但在晶格中，电子倾向于把正离子拉向自己，从而被一个晶格比平常略为密集的区域所包围。进入其邻近区域的另一个电子将被引向这个区域；看上去就好像它被第一个电子吸引着。打个比方：两个粒子紧挨着一同坐进床垫的同一个凹坑里，是能够占到便宜的。

计算这种力，最简单的办法是把它描述为一个电子发射“虚”声子、另一个电子吸收此声子的过程。考察图 200(a) 所定义的过程：处于状态 $\mathbf{k}$ 的电子发射一个声子，被散射到状态 $\mathbf{k}-\mathbf{K}$；处于状态 $\mathbf{k}'$ 的电子吸收这个声子，被散射到 $\mathbf{k}'+\mathbf{K}$。

%FIG{200}: {通过声子交换的电子–电子相互作用。}

在二级微扰论中，我们应当用矩阵元
\begin{equation*}
\langle \mathbf{k}-\mathbf{K},\,\mathbf{k}'+\mathbf{K}\,|\,V\,|\,\mathbf{k},\,\mathbf{k}'\rangle
=\frac{\mathcal{M}_{\mathbf{k},\,\mathbf{k}-\mathbf{K}}\,\mathcal{M}^{*}_{\mathbf{k}',\,\mathbf{k}'+\mathbf{K}}}{\mathcal{E}(\mathbf{k})-\mathcal{E}(\mathbf{k}-\mathbf{K})-\hbar\nu_{\mathbf{q}}}.
\tag{11.1}
\end{equation*}
来表示这一跃迁。在这个表达式中，$\mathcal{M}_{\mathbf{k},\mathbf{k}-\mathbf{K}}$ 代表电子–声子相互作用的矩阵元，与 §6.13 中的相同。能量分母代表由初态到中间态的能量变化，其中包括所产生的声子的能量 $\hbar\nu_{\mathbf{q}}$；对 N 过程有 $\mathbf{q}=\mathbf{K}$，不过电子–声子 U 过程也可以包括在内。

还有另一个过程对总散射作出贡献，如图 200(b) 所示：波矢为 $-\mathbf{q}$ 的声子先由 $\mathbf{k}'$ 发射出来，然后被 $\mathbf{k}$ 吸收。它给出一个与式 (11.1) 类似的表达式，只是分母稍有不同。把这两个过程合并起来，并假定电子–声子矩阵元实质上相同，我们就得到
\begin{equation*}
\langle \mathbf{k}-\mathbf{K},\,\mathbf{k}'+\mathbf{K}\,|\,V\,|\,\mathbf{k},\,\mathbf{k}'\rangle
=\frac{|\mathcal{M}_{\mathbf{k},\,\mathbf{k}-\mathbf{K}}|^{2}\,2\hbar\nu_{\mathbf{q}}}{\{\mathcal{E}(\mathbf{k})-\mathcal{E}(\mathbf{k}-\mathbf{K})\}^{2}-(\hbar\nu_{\mathbf{q}})^{2}}.
\tag{11.2}
\end{equation*}
它对电子所产生的效果，就好像两者之间存在一种直接的相互作用，其 Fourier 系数 $V(\mathbf{K})$ 由式 (11.2) 给出。它通常是正的；但对于
\begin{equation*}
|\mathcal{E}(\mathbf{k})-\mathcal{E}(\mathbf{k}-\mathbf{K})|<\hbar\nu_{\mathbf{q}}
\tag{11.3}
\end{equation*}
这一狭窄的能量范围，相互作用是负的，相应于一种吸引力。在式 (11.3) 的范围内，这个力具有 Fourier 分量
\begin{equation*}
V(\mathbf{K})\approx-\frac{2\,|\mathcal{M}_{\mathbf{k},\,\mathbf{k}-\mathbf{K}}|^{2}}{\hbar\nu_{\mathbf{q}}}.
\tag{11.4}
\end{equation*}
这或许算不上一种很强的力。我们是在极低的温度下工作，因此可以把声子发射过程看作由晶格的零点运动所诱发。利用式 (2.109)、(6.86) 与 (6.93)，我们得到
\begin{align*}
|\mathcal{M}_{\mathbf{k},\,\mathbf{k}-\mathbf{K}}|^{2}
&=|\mathcal{U}(\mathbf{K})|^{2}\,|\mathbf{K}\cdot\mathbf{U}_{\mathbf{q}}|^{2}\\
&=|\mathcal{U}(\mathbf{K})|^{2}\,\frac{\mathbf{K}^{2}\hbar}{2NM\nu_{q}},
\tag{11.5}
\end{align*}
式中 $\mathcal{U}(\mathbf{K})$ 是单个离子的屏蔽势的 Fourier 变换。如果只有 N 过程，并且声速由式 (6.84) 正确给出，那么
\begin{align*}
V(\mathbf{K}) &\approx-\frac{|\mathcal{U}(\mathbf{K})|^{2}\,\mathbf{q}^{2}}{NM\nu_{q}^{2}}\\
&=-\frac{|\mathcal{U}(\mathbf{K})|^{2}}{D_{0}s^{2}}=-\frac{3}{2}\,\frac{|\mathcal{U}(\mathbf{K})|^{2}}{NZ\mathcal{E}_{F}}.
\tag{11.6}
\end{align*}
这样，在 N 过程所允许的 $\mathbf{K}$ 值范围内，有效相互作用并不迅速减小。对更大的 $K$ 值，U 过程会使 $V(\mathbf{K})$ 有所增大，不过，正如 §6.13 和 §7.5 中所讨论的，电子 Bloch 态中平面波的混合避免了 $\mathbf{K}=\mathbf{g}$ 处的灾难。显然，这种吸引是短程的——直到两电子之间所能交换动量的最大值，它都具有很大的矩阵元。

当然，电子之间还存在着普通的库仑排斥。如同式 (5.32) 那样，相应的矩阵元应为
\begin{equation*}
V_{\mathrm{Coul.}}(\mathbf{K})=\frac{4\pi e^{2}}{K^{2}+\lambda^{2}},
\tag{11.7}
\end{equation*}
式中 $\lambda$ 是式 (5.22) 的屏蔽参量，或者是像式 (5.36) 那样的某种更复杂的表达式。我们把式 (11.4) 与式 (11.7) 之和取作对 \emph{Fröhlich 相互作用}（Fröhlich interaction）的估计。超导电性的判据就是它应为负值。用电子–声子相互作用来表述，这个判据并没有任何很简单的形式；不过有证据表明，Bardeen 公式 (6.93) 中“赝势部分”$\mathcal{V}'(K)$ 的强度才是关键因素。

上面的论证已经足够简单；但要想恰当地为它辩护，就需要考察电子–声子联合系统的能量中的各种贡献，其中包括 §6.11 中声速的重整化。应当注意，我们处理的只是像式 (11.3) 中那样处于很窄相对能量范围内的电子；在这个范围之外，晶格的运动太慢，跟不上电子的运动。但由于对净相互作用最重要的贡献来自相当大的 $K$ 值——相应于大的 $q$ 值——这个范围的量级为 $k\Theta$，这里 $\Theta$ 是 Debye 温度。既然超导电性只在极低温度下出现，这个能量范围就远比费米面（Fermi surface）上“热层”的厚度 $kT$ 宽广得多。因此，对于那些真正能够参与超导转变的电子来说，近似式 (11.4) 是合理的。

\section{Cooper 对}

为了领会电子之间交换声子所产生的吸引力的效果，我们现在来计算总散射振幅 $t$，即这样一个矩阵元：它的平方给出下述过程的概率——处于 $\mathbf{k}_1$ 和 $\mathbf{k}_2$ 两态的两个电子被散射到另外两个态 $\mathbf{k}_1'$ 和 $\mathbf{k}_2'$。如图 201(a) 那样，我们将假定这些态正好位于费米球（Fermi sphere）之外。

对散射最显而易见的贡献将是式 (11.2) 所定义的直接相互作用 $V(\mathbf{K})$。它对我们眼下的计算来说未免太复杂，所以让我们抽出其中的吸引部分，并假定它是一个常数。于是，我们的一级散射振幅将取为
\begin{equation*}
t_{1}=V(\mathbf{K})=V,
\tag{11.8}
\end{equation*}
其中 $V$ 是一个负的常数，与 $\mathbf{k}_1'$ 和 $\mathbf{k}_1$ 无关。但这个吸引势只对能量几乎相同的电子才成立；为了把条件 (11.3) 纳入理论，我们将约定：除非
\begin{equation*}
|\mathcal{E}(\mathbf{k}_{1})-\mathcal{E}(\mathbf{k}_{2})|<w,
\tag{11.9}
\end{equation*}
否则 $V(\mathbf{K})$ 为零，这里 $w$ 是一个常数。由式 (11.3) 显然可见，$w$ 应当具有 $k\Theta$ 的量级。

通常在散射问题中，近似式 (11.8) 已经足够。然而，在 $t$ 的微扰展开中还会有一些进一步的项，对应于电子从 $\mathbf{k}_1$ 出发、终于 $\mathbf{k}_1'$ 的逐次虚过程。例如，下一项可能就是图 201(b) 所示的过程：我们在初态与终态之间插入一个虚态，其中 $\mathbf{k}_1$ 先变到 $\mathbf{k}_1+\mathbf{K}$、$\mathbf{k}_2$ 先变到 $\mathbf{k}_2-\mathbf{K}$，然后再作第二次跃迁，变到 $\mathbf{k}_1'$ 和 $\mathbf{k}_2'$。

%FIG{201}: {(a) 电子的直接散射。(b) 二阶项。}

这种过程对散射振幅的贡献由通常的微扰表达式给出：
\begin{equation*}
t_{2}(\mathbf{K})=\frac{V^{2}}{\mathcal{E}(\mathbf{k}_{1})+\mathcal{E}(\mathbf{k}_{2})-\mathcal{E}(\mathbf{k}_{1}+\mathbf{K})-\mathcal{E}(\mathbf{k}_{2}-\mathbf{K})}.
\tag{11.10}
\end{equation*}
但这一项只在非常有限的条件下才存在。由于 Pauli 原理，中间态不得位于费米面以下；而且按式 (11.9)，每次跃迁的能量差必须小于 $w$。于是我们得到如下规则：
\begin{equation*}
\mathcal{E}_{F}<\mathcal{E}(\mathbf{k}_{1}+\mathbf{K})<\mathcal{E}(\mathbf{k}_{1})+w,\quad
\mathcal{E}_{F}<\mathcal{E}(\mathbf{k}_{2}-\mathbf{K})<\mathcal{E}(\mathbf{k}_{2})+w,
\tag{11.11}
\end{equation*}
等等；如果式 (11.10) 还想对 $t_2$ 有任何贡献的话，这些规则就必须得到满足。这些条件类似于 Kondo 效应（§10.6）中虚自旋翻转跃迁的选择定则；在同时计及直接过程和交换过程时，各个态的占据数并\emph{不}会从二阶项中消去。

这些规则限制性如此之强，以致通常我们可以忽略 $t_2$。但当 $\mathbf{k}_1$ 和 $\mathbf{k}_2$ 的方向近乎相反时，条件 (11.11) 的两式合二为一，对可能发生散射的 $\mathbf{K}$ 值范围也就不再施加那么苛刻的限制。事实上，此时 $t_2$ 可以变得相当大。

为说明这一点，方便的做法是从费米能级（Fermi level）起量度能量。我们对每个态定义
\begin{align*}
\xi(\mathbf{k}) &\equiv \mathcal{E}(\mathbf{k})-\mathcal{E}_{F}\\
&\approx \hbar v_{F}(k-k_{F})
\tag{11.12}
\end{align*}
（只要 $k$ 近似等于 $k_F$）。我们还引入两个具有能量量纲的新辅助变量
\begin{equation*}
\Delta=\xi(\mathbf{k}_{1})+\xi(\mathbf{k}_{2}),\qquad
\eta=\hbar v_{F}\,|\mathbf{k}_{1}+\mathbf{k}_{2}|
\tag{11.13}
\end{equation*}
它们量度的是我们这对电子的总能量和总动量。

为得到二阶微扰项，我们把式 (11.10) 在满足式 (11.11) 的所有 $\mathbf{K}$ 值上积分。利用式 (11.12) 和 (11.13)，这可以化成一个更简单的公式：
\begin{align*}
t_{2} &= V^{2}\int \frac{\mathrm{d}\mathbf{K}}{\mathcal{E}(\mathbf{k}_{1})+\mathcal{E}(\mathbf{k}_{2})-\mathcal{E}(\mathbf{k}_{1}+\mathbf{K})-\mathcal{E}(\mathbf{k}_{2}-\mathbf{K})}\\
&= \frac{2\pi V^{2}k_{F}^{2}}{8\pi^{3}\hbar v_{F}}\int_{0}^{w}\mathrm{d}\xi\int_{-1}^{1}\frac{\mathrm{d}(\cos\theta)}{\Delta+\eta\cos\theta-2\xi}.
\tag{11.14}
\end{align*}

%FIG{202}: {(a) 在 $\mathbf{k}_3$ 处产生一个空穴、随后又被 $\mathbf{k}_2$ 填充的过程。(b) 逐次的对散射。(c) 涉及空穴的对散射。}

当 $\Delta$ 和 $\eta$ 都远小于 $w$ 时，被积函数中出现一个奇点，它使积分呈对数行为（可参照式 (10.56)）：
\begin{equation*}
t_{2}\approx-\frac{V^{2}}{4\pi^{2}}\,\frac{k_{F}^{2}}{\hbar v_{F}}\ln\frac{2w}{\max{(\Delta,\eta)}}.
\tag{11.15}
\end{equation*}
通常这一项远小于 $t_1$，原本可以忽略。但如果 $\Delta$ 和 $\eta$ 恰好都非常小，式 (11.15) 就可能变得很大——这时我们或许不得不考虑 $t$ 的微扰级数中更进一步的项。

这样的项可能非常复杂，因为我们可以考虑如图 202(a) 那样的虚过程：一个来自费米面内部的电子与这对电子中的一个发生相互作用，留下一个空穴，而这个空穴最终在此后的某个过程中被填满。不过，这类过程在微扰级数中给出的贡献往往小得多。重要的是那些把图 201(b) 所示过程一再重复的情形——电子对不断地自我散射，如图 202(b) 那样。每一次这样的虚散射都引入一个类似于式 (11.15) 的因子，当 $\Delta$ 和 $\eta$ 很小时还包括那个大的对数积分。

为了得到正确的结果，我们还必须把这样的过程包括进来：一对电子从费米海（Fermi sea）中被散射出去，留下一对空穴，它们可以接收 $\mathbf{k}_1$ 和 $\mathbf{k}_2$ 处的电子（图 202(c)）。这一过程对 $t_2$ 给出与式 (11.15) 相同的贡献。把所有这些项加起来，我们看到散射振幅必定表现为无穷级数
\begin{equation*}
t\approx V+\gamma V^{2}+\gamma^{2}V^{3}+\ldots,
\tag{11.16}
\end{equation*}
其中
\begin{equation*}
\gamma=-\frac{1}{2\pi^{2}}\,\frac{k_{F}^{2}}{\hbar v_{F}}\ln\frac{2w}{\max{(\Delta,\eta)}}.
\tag{11.17}
\end{equation*}
但级数 (11.16) 可以求和：
\begin{equation*}
t\approx\frac{V}{1-\gamma V}.
\tag{11.18}
\end{equation*}
这样，散射振幅的行为就好像它含有一个极点。这是众所周知的一个征兆，表明自由粒子彼此散射的假设已经失效；它正是两个电子形成\emph{束缚态}（bound state）的标志。假设我们取 $\eta=0$，使得 $\mathbf{k}_1=-\mathbf{k}_2$。极点出现时的能量由下式确定：
\begin{align*}
1&=\gamma V\\
&=-\frac{V}{2\pi^{2}}\,\frac{k_{F}^{2}}{\hbar v_{F}}\ln\frac{2w}{\Delta}.
\tag{11.19}
\end{align*}
（译注：原书式 (11.19) 第二行分母印作 $v_F$，漏排 $\hbar$；由式 (11.17) 显见应有 $\hbar v_F$，译文已补正。）

换句话说，当这对电子的总能量在费米能级以下等于
\begin{equation*}
\Delta_{0}=2w\exp\left(-\frac{2\pi^{2}\hbar v_{F}}{|V|\,k_{F}^{2}}\right)
\tag{11.20}
\end{equation*}
时，它们就倾向于彼此束缚而形成一个“准分子”（quasi-molecule）。这样的状态称为 \emph{Cooper 对}（Cooper pair）。

这类态有可能存在，但这并没有解决超导电性的问题。我们可以这样论证：金属中所有的电子都倾向于结成这样的对——$\mathbf{k}$ 与 $-\mathbf{k}$ 相配，而且这种由 Bose–Einstein 型“准分子”组成的气体随后便能发生 Bose–Einstein 凝聚。然而这些电子对并不稳定；作为对超导态的描述，这种说法有启发性，却缺乏说服力。

最重要的结果是关于电子对能量的公式 (11.20)。我们可以采用自由电子模型，它给出
\begin{align*}
\frac{k_{F}^{2}}{2\pi^{2}\hbar v_{F}} &= \frac{1}{2}\,\mathcal{N}(\mathcal{E}_{F})\\
&= \frac{3}{4}\,\frac{ZN}{\mathcal{E}_{F}}
\tag{11.21}
\end{align*}
